\begin{ledgroupsized}[r]{120mm}
\footnotesize
\pstart
\noindent\textbf{Überlieferung:}
\pend
\end{ledgroupsized}
\begin{ledgroupsized}[r]{114mm}
\footnotesize
\pstart \parindent -6mm
\makebox[6mm][l]{\textit{L}}%
Aufzeichnung:
LH~XXXVII~5 Bl.~198. 
Ein Blatt~8\textsuperscript{o};
\href{\Scan{LH_37_05_198wz}}%
{Wasserzeichenfragment}
am Blattrand;
Papiererhaltungsmaßnahmen;
alle Ränder beschnitten.
Eine Seite auf Bl.~198~r\textsuperscript{o}; 
auf Bl.~198~v\textsuperscript{o} Bruchstück einer Adresse (Hand eines unbekannten Korrespondenten): 
\textit{A} \lbrack/\rbrack\ 
\textit{Monsi}\textlangle\textit{eur}\textrangle\ 
\textlangle\textendash\textrangle\ \lbrack/\rbrack\ 
\textit{Conseille}\textlangle\textit{r}\textrangle\ \textlangle\textendash\textrangle\  \lbrack/\rbrack\ 
\textit{l'Elec}\textlangle\textit{teur}\textrangle\ \textlangle\textendash\textrangle.
\pend
\end{ledgroupsized}
\vspace{5mm}
\begin{ledgroup}
\footnotesize
\pstart
\noindent%
\textbf{Datierungsgründe:}
Die Aufzeichnung verweist auf Leibnizens 
\cite{01502}Brief an D.~Papin\indexN{\textls{Papin} (Papinus), Denis 1647\textendash1713} vom 8.\ (18.) November 1697 (\textit{LSB} III,~7 N.\;156)
und auf die darin geführte Diskussion über die Rolle der Elastizität beim Stoß, am Beispiel einer zwischen den stoßenden Körpern befindlichen Feder.
Dieses Beispiel wird zunächst von Papin 
am 2. (12.) Juli 1696
vorgeschlagen
(\cite{02070}\textit{LSB} III,~7 N.\;2, hier S.\;20),
von Leibniz  am 16. (26.) Juli 1696 aufgegriffen (\cite{02071}\textit{LSB} III,~7 N.\;8)
und von beiden Kontrahenten in fast allen Briefen der folgenden Monate rege diskutiert. 
Nach dem in N.\;\ref{60360} erwähnten Schreiben vom 8./18. November 1697
wird der besondere Stoßfall in Leibnizens Brief vom 2. (12.) Dezember 1697 (\cite{02072}N.\;163) sehr ausführlich behandelt und im Brief vom 16. (26.) Januar 1698 (\cite{02073}N.\;177) ein letztes Mal erwähnt.
Die vorgeschlagene Datierung beruht auf der Annahme einer Entstehung von N.\;\ref{60360} im Rahmen der geschilderten Diskussion bzw. kurze Zeit später.
\pend
\end{ledgroup}
\selectlanguage{latin}
\vspace{8mm}
\pstart%
\normalsize%
\noindent%
\href{\Scan{LH_37_05_198r}}%
{\lbrack198~r\textsuperscript{o}\rbrack}\
Corpus\indexS{corpus}
\edtext{corpori ingruit,}{\lemma{corpori}\Bfootnote{%
\textit{(1)}~\textls{ingruit} %
\textit{(2)}~ingruit,~\textit{L}}} 
quamdiu corpora sibi appropinquant, sive
\edtext{ea}{\lemma{}\Bfootnote{ea \textit{erg.~L}}} 
\indexS{corpora in easdem partes tendentia}tendant in easdem partes, 
\indexS{corpora in contrarias partes tendentia}sive in contrarias. 
Sed si in contrarias tendant, utrumque alteri ingruere dicetur; 
sin tendant in easdem partes, id demum alteri ingruet, quod movetur celerius, seu quod persequitur alterum.%
\pend
\pstart
Itaque
\edtext{generaliter, \indexS{corpus ingruens}\textls{ingruit},}{\lemma{generaliter}\Bfootnote{%
\textit{(1)}~\textls{ingruet}, %
\textit{(2)}~\textls{ingruit},~\textit{L}}} 
in quo est ratio appropinquationis\indexS{appropinquatio}; cujus scilicet is est status, ut altero
\edtext{licet}{\lemma{}\Bfootnote{licet \textit{erg.~L}}} 
ad \indexS{quies}quietem reducto, tamen appropinquatio\indexS{appropinquatio} perseveret. 
\edtext{Usus hujus notionis generalis a me deprehensus est cum de Elastro a corporibus concurrentibus tenso\indexS{elastrum a corporibus concurrentibus tenso} ratiocinarer.}{\lemma{cum de Elastro \lbrack...\rbrack\ ratiocinarer}\Cfootnote{%
Siehe Leibnizens \cite{01502}Brief an D. Papin vom 8. (18.) November 1697, \textit{LSB} III,~7 N.\;156.%
}}
\edtext{Reperi nempe tamdiu Elastrum\indexS{elastrum} 
\edtext{\lbrack magis\rbrack}{%
\lemma{}%
\Bfootnote{%
magis %
\textit{erg. Hrsg.}}}
magisque intendi quamdiu superest ingruentia\indexS{ingruentia}, 
seu donec corpus quod adhuc ingruit\indexS{corpus ingruens} reducatur ad quietem\indexS{quies}.}{\lemma{Reperi \lbrack...\rbrack\ quietem}\Cfootnote{%
Leibniz an Papin, 16. (26.) Juli 1696, \textit{LSB} III,~7 N.\;8, bes. S.\;34.%
}}
Neque enim ad contrarium motum deducitur donec ad quietem\indexS{quies} 
\ledpb%
reducatur.
\edtext{A solo etiam ingruente corpore\indexS{corpus ingruens}
Elastrum\indexS{elastrum} urgetur. 
Etsi revera eadem sit elastri\indexS{elastrum} tensio quodcunque sit urgens 
vel ingruens, modo eadem appropinquatio\indexS{appropinquatio} servetur.}{\lemma{}\Bfootnote{A solo \lbrack...\rbrack\ revera %
\textit{(1)}~nihil %
\textit{(a)}~refert %
\textit{(b)}~referat %
\textit{(aa)}~quoad %
\textit{(bb)}~quodnam %
\textit{(2)}~eadem sit elastri tensio quodcunque sit \lbrack...\rbrack\ servetur. \textit{erg.~L}}} 
Videantur
\edtext{de his}{\lemma{}\Bfootnote{de his \textit{erg.~L}}} 
\edtext{ea quae commentatus sum Novembr.\ 1697.\ cum
Dni Papini\indexN{\textls{Papin} (Papinus), Denis 1647\textendash1713} literis responderem.}{\lemma{ea \lbrack...\rbrack\ responderem}\Cfootnote{\cite{01502}Brief an D. Papin vom 8. (18.) November 1697, \textit{LSB} III,~7 N.\;156, bes. S.\;633\textendash635.}}
\pend 
